\subsection{Kod do fitowania wyników z eksperymentu}
\begin{python}
import pandas as pd
from zipfile import Path as ZipPath
import matplotlib.pyplot as plt
from pathlib import Path
import pandas as pd
import numpy as np
from scipy.optimize import curve_fit
import pandas as pd



plt.rcParams.update({
	"font.family": "serif",  # use serif/main font for text elements
	"text.usetex": True,     # use inline math for ticks
	"pgf.rcfonts": False     # don't setup fonts from rc parameters
})
fps = {x.stem: x for x in ZipPath('Dopasowanie.zip', at='/').iterdir()}
fps     
for k,v in fps.items():
print(k, v.exists())
%%time

dfs = {}
for k,v in fps.items():
with v.open() as f:
dfs[k] = pd.read_csv(f, parse_dates=['UTC'], index_col='UTC')
df = pd.concat(
[
dfs["IRL-2.1"][["Flux, cps"]].rename(columns={"Flux, cps": "IRL-2.1"}),
dfs["IRL-2.2"][["Flux, cps"]].rename(columns={"Flux, cps": "IRL-2.2"}),
dfs["trpp"]
],
axis=1
).sort_index()
print(df.columns.tolist())
print(dfs.keys())

%%time
fig, (ax1, ax2) = plt.subplots(nrows=2, figsize=(12,12), sharex=True)


df.plot(
ax=ax1,
y=["IRL-2.1", "IRL-2.2"],
logy=True,
grid=True,
marker=".",
)

# DÓŁ: TRPP
df.plot(
ax=ax2,
y=["PAR", "PK2", "PK6", "PB1"],
grid=True,
marker=".",
)

ax1.legend()
ax2.legend()
plt.tight_layout()
plt.show()


ax1.legend()
intervals = {
	"zrzuty": slice("2025-12-08 19:55:30", "2025-12-08 21:30:30"),
	"zrzut01": slice("2025-12-08 19:55:30", "2025-12-08 20:03:30"),
	"zrzut02": slice("2025-12-08 20:07:30", "2025-12-08 20:11:30"),
	"zrzut03": slice("2025-12-08 20:15:30", "2025-12-08 20:19:30"),
	"zrzut04": slice("2025-12-08 20:23:30", "2025-12-08 20:27:30"),
	"zrzut05": slice("2025-12-08 20:31:30", "2025-12-08 20:35:30"),
	"zrzut06": slice("2025-12-08 20:39:30", "2025-12-08 20:43:30"),
	"zrzut07": slice("2025-12-08 20:47:30", "2025-12-08 20:51:30"),
	"zrzut08": slice("2025-12-08 20:55:30", "2025-12-08 20:59:30"),
	"zrzut09": slice("2025-12-08 21:00:30", "2025-12-08 21:07:30"),
	"zrzut10": slice("2025-12-08 21:08:30", "2025-12-08 21:12:30"),
	"zrzut11": slice("2025-12-08 21:12:30", "2025-12-08 21:18:30"),
}

[(k, v.start, v.stop) for k, v in intervals.items()]
import matplotlib.pyplot as plt
import matplotlib.dates as mdates

for name, sl in intervals.items():

df_cut = df.loc[sl]
if name == "zrzuty":
continue
if df_cut.empty:
print(f"{name}: pusty zakres")
continue

fig, (ax1, ax2) = plt.subplots(
nrows=2,
figsize=(16, 12),
sharex=True
)


df_cut.plot(
ax=ax1,
y=["IRL-2.1"],
logy=True,
grid=True,
marker=".",
)


df_cut.plot(
ax=ax2,
y=["PK6"],
grid=True,
marker=".",
)

# format osi X
ax1.set_title(f"{name} | {sl.start} → {sl.stop}")
ax2.set_xlabel("UTC")

plt.tight_layout()
plt.show()

df_cut = df.loc[intervals["zrzut01"]]

fig, (ax1, ax2) = plt.subplots(nrows=2, figsize=(12,12), sharex=True)

df_cut.plot(
ax=ax1,
y=["IRL-2.1"],
logy=True,
grid=True,
marker=".",
)

df_cut.plot(
ax=ax2,
y=["PK6"],
grid=True,
marker=".",
)

plt.tight_layout()
plt.show()
cuts = {
	"zrzut01": ("2025-12-08 19:57:30", "2025-12-08 20:01:30"),
	"zrzut02": ("2025-12-08 20:08:40", "2025-12-08 20:09:45"),
	"zrzut03": ("2025-12-08 20:17:15", "2025-12-08 20:18:05"),
	"zrzut04": ("2025-12-08 20:25:05", "2025-12-08 20:26:10"),
	"zrzut05": ("2025-12-08 20:32:48", "2025-12-08 20:33:43"),
	"zrzut06": ("2025-12-08 20:41:08", "2025-12-08 20:41:53"),
	"zrzut07": ("2025-12-08 20:49:13", "2025-12-08 20:50:08"),
	"zrzut08": ("2025-12-08 20:56:33", "2025-12-08 20:57:33"),
	"zrzut09": ("2025-12-08 21:02:50", "2025-12-08 21:03:40"),
	"zrzut10": ("2025-12-08 21:09:22", "2025-12-08 21:10:12"),
	"zrzut11": ("2025-12-08 21:14:59", "2025-12-08 21:15:34"),
}

for name, (t0, t1) in cuts.items():
maly_df = df.loc[t0:t1]

t = (maly_df.index - maly_df.index[0]).total_seconds()
y = maly_df["IRL-2.1"]
t
from scipy.optimize import curve_fit
import numpy as np
import matplotlib.pyplot as plt

b_i   = np.array([2.43e-4, 1.363e-3, 1.203e-3, 2.605e-3, 8.19e-4, 1.67e-4])
lambda_i = np.array([1.27e-2, 3.17e-2, 1.15e-1, 3.11e-1, 1.40,   3.87   ])  # [1/s]
B = b_i.sum()
L = 146e-6

def exp_fit(t, a, T):
return np.exp(t/T) * a   

def rho_T(T_hat, L,B, b_i, lambda_i):
suma = np.sum((b_i) / ((1 /T_hat)+lambda_i))
nie_suma = 1/(T_hat*B)
rho = nie_suma * (L+ suma)
return rho

def sigma_rho_T(sigma_T, L,B, b_i, lambda_i, T_hat):
suma = np.sum((b_i * lambda_i) / (B * (lambda_i*T_hat + 1)**2))
nie_suma = L/((T_hat**2)*B)
sigma_rho = (nie_suma + suma)*sigma_T
return sigma_rho

def pm(val, err, vfmt="{:.2f}", efmt="{:.2f}"):
return f"{vfmt.format(val)} ± {efmt.format(err)}"

wyniki = []

fig, ax = plt.subplots(figsize=(20,12))

for name, (t0, t1) in cuts.items():
maly_df = df.loc[t0:t1, ["IRL-2.1"]].dropna()
if maly_df.empty:
print(name, "-> pusty wycinek")
continue

t = (maly_df.index - maly_df.index[0]).total_seconds().to_numpy()
y = maly_df["IRL-2.1"].to_numpy()

if not (np.isfinite(t).all() and np.isfinite(y).all()):
print(name, "-> NaN/inf w t lub y")
continue


try:
popt, pcov = curve_fit(exp_fit, t, y, p0=[y[0], 60], maxfev=10000)
except Exception as e:
print(name, "-> curve_fit error:", e)
continue

sigma_a, sigma_T = np.sqrt(np.diag(pcov))
a_hat, T_hat = popt

r  = rho_T(T_hat, L, B, b_i, lambda_i)
dr = sigma_rho_T(sigma_T, L, B, b_i, lambda_i, T_hat)

wyniki.append([name, T_hat, sigma_T, r, dr])

ax.plot(t, y, marker='.', ls='None', ms=2, alpha=0.6, label=f"{name} data")
ax.plot(t, exp_fit(t, *popt), lw=2, label=f"{name} fit")

# format wykresu
ax.set_yscale('log')
ax.grid(True, which='both', alpha=0.3)
ax.legend(ncol=2, fontsize=18)
ax.set_xlabel("Czas od początku zrzutu s")
ax.set_ylabel("IRL-2.1 Flux, cps")
plt.show()

# --- tabela wyników ---
import pandas as pd
tab = pd.DataFrame(wyniki, columns=["zrzut", "T_hat [s]", "sigma_T [s]", "rho", "delta_rho"])
print(tab)

tab["T ± ΔT"]   = [pm(T, dT, vfmt="{:.3f}", efmt="{:.3f}") for T, dT in zip(tab["T_hat [s]"], tab["sigma_T [s]"])]
tab["ρ ± Δρ"]   = [pm(r, dr, vfmt="{:.6f}", efmt="{:.6f}") for r, dr in zip(tab["rho"], tab["delta_rho"])]
print(tab[["zrzut", "T ± ΔT", "ρ ± Δρ"]])


outdir = "/home/piegus/ncbj/Inzynierka/figures/"
# Save plot
plt.savefig(f"{outdir}/zrzuty.pgf", format='pgf', bbox_inches='tight')

plt.show()

\end{python}
\subsection{Kod do wyznaczenia miejsc zerowych }
\begin{python}

from __future__ import annotations
import numpy as np
import pandas as pd
import matplotlib.pyplot as plt
from scipy.integrate import solve_ivp
from numpy.polynomial import Polynomial
from ipywidgets import interact, Dropdown
import ipywidgets as W
import matplotlib.ticker as ticker
from math import sqrt
import matplotlib as mpl

# 0b.  USTAWIENIA MATPLOTLIB
mpl.rcParams.update({
	'font.family': 'serif',
	'font.serif': ['DejaVu Serif', 'Libertinus Serif', 'CMU Serif', 'Nimbus Roman'],
	'mathtext.fontset': 'cm',     # ładny wygląd symboli i 10^n bez LaTeX
	'axes.unicode_minus': False,   
	
	'font.size': 20,
	'axes.labelsize': 20,
	'legend.fontsize': 20,
	'xtick.labelsize': 14,
	'ytick.labelsize': 14,
	
	'figure.dpi': 140,
	'savefig.dpi': 300,
	'figure.constrained_layout.use': True,  #
})


beta_i   = np.array([2.43e-4, 1.363e-3, 1.203e-3,
2.605e-3, 8.19e-4, 1.67e-4])
lambda_i = np.array([1.27e-2, 3.17e-2, 1.15e-1,
3.11e-1, 1.40,   3.87   ])

beta          = beta_i.sum()         # 0.006552 (1 $)
Lambda_prompt = 150e-6               # 150 µs
t_analysis    = 0.0


def rho_factory(rho_max: float):
"""Zwraca rho(t) = rho_max [$] jako skok jednostkowy w chwili t=0."""
def rho(t: float) -> float:
return rho_max if t >= 0 else 0.0
return rho

rho_max_1 = 1.5
rho_max_2 = 1.0
rho_max_3 = 0.5
rho_max_4 = 0.1
rho_max_10 = 0.2
rho1 = rho_factory(rho_max_1)
rho2 = rho_factory(rho_max_2)
rho3 = rho_factory(rho_max_3)


def kinetics_matrix(t: float, rho_func) -> np.ndarray:
r = rho_func(t)  # r w $, czyli ρ = r * β
A = np.zeros((7, 7))
A[0, 0]  = (r*beta - beta) / Lambda_prompt   # (ρ - β)/Λ (dla ρ w $)
A[0, 1:] = lambda_i
A[1:, 0] = beta_i / Lambda_prompt
A[1:, 1:] = -np.diag(lambda_i)
return A

def I(n: int = 7) -> np.ndarray:
return np.eye(n)

def Matrix_alpha(t: float, alpha: float, rho_func) -> np.ndarray:
return kinetics_matrix(t, rho_func) - alpha * I()

def p_alpha(alpha, t, rho_func):
"""Wielomian charakterystyczny p(α) = det(A - αI)."""
alpha = np.atleast_1d(alpha)
return np.array([np.linalg.det(Matrix_alpha(t, a, rho_func)) for a in alpha])
\end{python}
\subsection{Przebieg mocy reaktora w czasie}
\begin{python}
	import numpy as np
	import matplotlib.pyplot as plt
	import matplotlib.ticker as ticker
	
	t_step = 50  # chwila skoku ρ
	

	T_EXP_SLOPE = {0.5: 1.85, 1.0: 0.832, 1.5: 0.182}
	
	def plot_case_sum(
	rho_val, t_range, color,
	show_step=True,
	show_dom_segment=True,
	segment_rel_len=0.02,
	t_mark_rel=None,            #
	):
	"""
	n(t)/n0 po skoku ρ w t_step: pełna suma modów
	"""

	t_start, t_end = t_range
	t = np.linspace(t_start, t_end, max(2, int((t_end - t_start)*400)))
	

	n0 = 1.0
	C0 = beta_i / (Lambda_prompt * lambda_i) * n0
	y0_state = np.concatenate(([n0], C0))
	
	rho_fun = RHO_FUN[rho_val] if 'RHO_FUN' in globals() else rho_factory(rho_val)
	A_post = kinetics_matrix(0.0, rho_fun)
	alpha_k, V = np.linalg.eig(A_post)
	a_k = np.linalg.solve(V, y0_state)
	N_k = a_k * V[0, :]   # wkłady neutronowe
	
	ak = np.real(alpha_k)
	Nk = np.real(N_k)
	

	y = np.ones_like(t, dtype=float)
	m = t >= t_step
	if np.any(m):
	tau = t[m] - t_step
	exp_arg = np.outer(tau, ak)
	exp_arg = np.clip(exp_arg, -700.0, 700.0)   # anty-over/underflow
	n_full  = (Nk * np.exp(exp_arg)).sum(axis=1)
	y[m] = np.clip(n_full, 1e-300, 1e300)
	

	y[y < 1.0] = np.nan
	

	fig, ax = plt.subplots(figsize=(12, 8))
	ax.plot(t, y, color=color, lw=2, label=fr"$\rho = {rho_val:.2f}\ \$$")
	
	ax.set_yscale("log")
	ax.yaxis.set_major_locator(ticker.LogLocator(base=10))
	ax.yaxis.set_major_formatter(
	ticker.FuncFormatter(lambda v, pos: ('1' if np.isclose(v, 1.0) else r'$10^{%d}$' % int(np.log10(v))) if (v > 0) else '')
	)
	ax.yaxis.set_minor_formatter(ticker.NullFormatter())
	
	y_pos = y[(y > 0) & np.isfinite(y)]
	if y_pos.size >= 2:
	ymax = min(float(y_pos.max())*1.05, 1e3)
	top = 10**np.ceil(np.log10(ymax))
	ax.set_ylim(0.85, min(1e3, top))
	
	else:
	ax.set_ylim(0.01, 10.0)
	
	if show_step:
	ax.axvline(t_step, color="0.3", ls="--", lw=1, label="skok ρ")
	

	mask_post = np.isfinite(y) & (y > 0) & (t >= t_step)
	t_post = t[mask_post]; y_post = y[mask_post]
	if t_post.size >= 5:
	k_tail = max(5, int(np.ceil(0.20 * t_post.size)))
	t_fit  = t_post[-k_tail:]; y_fit = y_post[-k_tail:]
	slope, intercept = np.polyfit(t_fit, np.log(y_fit), 1)  # ln(y)=slope*t + intercept
	
	t_line = np.linspace(t_start, t_end, 400)
	y_line = np.exp(intercept + slope * t_line)
	ax.plot(t_line, y_line,
	linestyle=(0, (9, 5)), color="black", lw=2.0, alpha=0.9, zorder=9,
	label=" rozwiązanie asymptotyczne")
	

	y_line_pos = y_line[(y_line > 0) & np.isfinite(y_line)]
	if y_line_pos.size:
	y_all_max = max(float(np.nanmax(y_pos)) if y_pos.size else 1.0,
	float(np.nanmax(y_line_pos)))
	ax.set_ylim(0.85, min(1e3, y_all_max * 1.05))
	

	tx = t_step + (t_mark_rel if (t_mark_rel is not None) else 0.8*(t_end - t_start))
	tx = np.clip(tx, t_start, t_end)
	yx = np.exp(intercept + slope * tx)
	
	# kąt linii w pikselach (poprawny dla osi log)F
	angle_offset_deg = -3 # możesz doprecyzować np. +2.0 / -1.5
	dt_px = 10
	trans = ax.transData
	X, Y = trans.transform((tx, yx))
	x2_data, _ = trans.inverted().transform((X + dt_px, Y))
	y2_data = np.exp(intercept + slope * x2_data)
	X2, Y2 = trans.transform((x2_data, y2_data))
	angle = np.degrees(np.arctan2(Y2 - Y, X2 - X)) + angle_offset_deg
	
	ax.annotate(
	r"$n(t)\ \propto\ e^{\dfrac{t}{T}}$",
	xy=(tx, yx), xycoords='data',
	xytext=(0, 14), textcoords='offset points',  # „nad” linią
	rotation=angle, rotation_mode='anchor',
	ha='left', va='bottom',
	size=20,
	zorder=10, clip_on=False
	)
	
	ax.set_xlim(t_start, t_end)
	ax.set_xlabel("Czas, s")
	ax.set_ylabel(r"$n(t)/n_0$")
	ax.grid(True, which="both", ls=":", alpha=0.7)
	ax.legend(loc="lower right")
	plt.tight_layout()
	plt.show()
	

	plot_case_sum(0.5, (47.5, 53), "purple", t_mark_rel=T_EXP_SLOPE[0.5])
	plot_case_sum(1.0, (48, 53),    "blue",  t_mark_rel=T_EXP_SLOPE[1.0])
	plot_case_sum(1.5, (49.8, 50.6),"red",   t_mark_rel=T_EXP_SLOPE[1.5])
	
\end{python}

